\textbf{Servicio de reabastecimiento declarado.} La oferta propone Enex Express como prestador principal de entrega de diésel en terreno y Copec Empresas como alternativa de suministro industrial. Sus servicios públicos respaldan el tipo de prestación; la cobertura específica de Bahía Panitao y los plazos siguientes corresponden a condiciones contractuales propuestas que Synaptix debe formalizar antes de la puesta en servicio \parencites{p6-enex-express}{p6-copec-industrial}.

El servicio cubrirá los dos estanques propios de generación, con atención de pedidos y coordinación de despacho las veinticuatro horas. Synaptix supervisa niveles, activa el pedido, confirma despacho y coordina acceso y recepción con el CLIENTE; un operador competente ejecuta descarga y control de calidad. Cada estanque tendrá al menos 1.190 L útiles. Con consumo de cálculo de 41,3 L/h, esa reserva representa $1.190/41,3=28,81$ horas; sostiene veinticuatro horas sin depender de una entrega durante ese intervalo. El estanque estándar de 394 L del grupo no satisface esta autonomía \parencite[pp.~2--3]{p7-fg-p200-ficha}.

La reposición se activa al llegar a 450 L útiles o cuando la proyección de demanda anticipa ese umbral. El plazo propuesto es de ocho horas desde el pedido, con confirmación en una hora. Si no se confirma, Synaptix activa el prestador alternativo al término de esa hora y conserva el mismo vencimiento de ocho horas contado desde el pedido inicial. A máxima demanda, $450-8(41,3)=119,6$ L proporcionan 2,90 horas adicionales después del vencimiento; la entrega debe restituir al menos 1.190 L útiles por estanque. Nivel, pedido, confirmación, volumen recibido y calidad quedan registrados.

La contratación especificará diésel compatible con el motor, volumen por entrega, acceso terrestre a ambos puntos de descarga, vehículo y personal autorizados, contención, seguridad, limpieza, tratamiento de agua y conciliación de volumen. Antes de una marejada anunciada se completa la reserva y se confirma la ruta alternativa. Si se prevé un bloqueo de acceso superior a la reserva efectiva, se amplía el almacenamiento autorizado o se prepara una provisión segura previa; compartir combustible del surtidor comercial no es la contingencia prevista. El servicio contratado cubre la vigencia completa de operación y se formaliza antes de habilitación, con acceso ordinario y contingencia. La contratación y la prueba de despacho deben acreditar los plazos y la factibilidad local antes de habilitación, sin atribuir al prestador un acuerdo ya firmado.
